\documentclass[11pt]{article}

\usepackage[a4paper,margin=1in]{geometry}
\usepackage{graphicx}
\usepackage{amsmath}
\usepackage{booktabs}
\usepackage{hyperref}

\title{Predicting Patient Diagnosis Using Machine Learning}
\author{Laura Valentina Sierra Pena\and Ainur Smailova \and Saif Ur Rehman}
\date{\today}

\begin{document}

\maketitle

\section{Problem Formulation}

The objective of this study is to predict patient diagnostic outcomes using demographic, behavioural, and clinical variables. The target variable, \textit{diagnosis}, contains three categories: \textit{Depression}, \textit{CI}, and \textit{Both}. Therefore, the task is formulated as a supervised multiclass classification problem.

Machine learning models are trained to learn relationships between patient characteristics and diagnostic outcomes. The goal is to develop a predictive model capable of accurately classifying unseen patient cases.

\section{Data Management}

\subsection{Exploratory Data Analysis}

Initial exploration was conducted to understand the dataset structure and detect potential data quality issues. The dataset contains 670 observations and 19 variables representing demographic information, behavioural patterns, and clinical measurements.

Missing value analysis confirmed that no missing values were present in the dataset, meaning no imputation strategies were required.

Multiple visualisation techniques were used to analyse the data:

\begin{itemize}
\item Histograms were used to examine the distribution of numerical variables.
\item Violin plots were used to assess distribution shape and spread.
\item Boxplots were used to detect potential outliers.
\end{itemize}

Although several extreme values were observed, these values were within clinically valid ranges specified in the dataset description and were therefore retained.

\subsection{Feature Screening}

A correlation heatmap was used to detect multicollinearity among numerical variables. A very strong correlation was observed between \textit{sleep\_quality\_index} and \textit{deep\_sleep\_quality\_index}. Since highly correlated variables introduce redundancy, the variable \textit{deep\_sleep\_quality\_index} was removed.

Categorical relationships were analysed using Cramér’s V statistics. This analysis revealed negligible associations among most categorical predictors. However, a perfect association was identified between the variable \textit{pem\_present} and the diagnosis variable. This indicates target leakage, meaning the feature contains direct information about the outcome. Consequently, \textit{pem\_present} was removed from the dataset.

The variable \textit{id} was also removed as it represents a unique identifier and does not provide predictive information.

\subsection{Feature Encoding}

Machine learning algorithms require numerical input; therefore categorical variables were encoded before model training.

Nominal variables such as \textit{gender} and \textit{work\_status} were encoded using one-hot encoding. Ordinal variables including \textit{social\_activity\_level} and \textit{exercise\_frequency} were encoded using ordinal encoding to preserve their natural ordering. Binary variables were encoded as indicator values.

To prevent data leakage, preprocessing was implemented within model pipelines and applied only during training folds in cross-validation.

\subsection{Experimental Protocol}

To evaluate model performance reliably, the dataset was divided into training and testing subsets using an 80/20 stratified split. Stratification ensures that class proportions remain consistent across both sets.

Model selection and hyperparameter tuning were performed using stratified 5-fold cross-validation on the training set. Macro F1-score was used as the primary evaluation metric because it gives equal importance to all classes.

\section{Model Training}

Three machine learning models discussed during lectures were selected for comparison:

\begin{itemize}
\item Logistic Regression
\item k-Nearest Neighbours (KNN)
\item Random Forest
\end{itemize}

These models represent different learning paradigms: linear classification, instance-based learning, and ensemble tree-based learning.

Hyperparameter tuning was conducted using Grid Search. For Logistic Regression, the regularisation parameter $C$ was tuned. For KNN, the number of neighbours and weighting scheme were optimised. For Random Forest, the number of trees and maximum tree depth were explored.

Cross-validation results showed that Random Forest achieved the highest Macro F1-score, significantly outperforming Logistic Regression and KNN. This indicates that nonlinear relationships and feature interactions are important for predicting patient diagnosis.

Therefore, Random Forest was selected as the final model.

\section{Model Evaluation}

The selected Random Forest model was evaluated on the unseen test set.

The model achieved the following performance:

\begin{itemize}
\item Test Accuracy: 0.963
\item Macro F1-score: 0.954
\end{itemize}

Class-wise evaluation showed that the CI class was perfectly classified. The Depression class also achieved high precision and recall. The Both category exhibited slightly lower recall, with a small number of cases misclassified as Depression.

The confusion matrix indicates that classification errors are minimal and primarily occur between the Both and Depression categories. This may reflect overlapping clinical characteristics between these groups.

Overall, the Random Forest model demonstrates strong predictive capability and generalises well to unseen data.

\section{Conclusion}

This study applied several machine learning techniques to predict patient diagnostic outcomes based on demographic and clinical variables. Through careful exploratory data analysis, feature screening, and rigorous model evaluation, Random Forest was identified as the most effective model.

The model achieved high accuracy and balanced performance across all classes, demonstrating that ensemble tree-based methods are well suited for structured healthcare datasets.

Future work could explore additional ensemble methods, larger datasets, and more advanced feature engineering techniques to further improve predictive performance.

\end{document}